\documentclass[11pt]{article}
\usepackage[letterpaper,margin=0.85in]{geometry}\usepackage[T1]{fontenc}\usepackage[utf8]{inputenc}\usepackage{amsmath,amssymb,enumitem,booktabs,array,xcolor,fancyhdr,tabularx}
\setlist{itemsep=2pt,topsep=3pt}\pagestyle{fancy}\fancyhf{}\renewcommand{\headrulewidth}{0.3pt}
\fancyhead[L]{\small Free-Body Diagrams and Net Force}\fancyhead[R]{\small Teacher key}\fancyfoot[C]{\small\thepage}
\definecolor{ans}{RGB}{28,125,74}\newcommand{\A}[1]{\textcolor{ans}{\textbf{#1}}}
\begin{document}

\section*{Free-Body Diagrams and Net Force --- teacher key}
\textbf{Question:} Which forces act on an object, and which way do they add up?\par
\subsection*{Outline and pacing (42 min)}\begin{tabular}{@{}llll@{}}\toprule Tab & Priority & Minutes & If behind\\\midrule
start & 1 & 2 & never trimmed\\
warmup & 1 & 6 & never trimmed\\
predict & 1 & 3 & never trimmed\\
explore & 3 & 4 & skipped first\\
build & 2 & 10 & shrinks to 6 of 12\\
practice & 1 & 5 & never trimmed\\
check & 1 & 3 & never trimmed\\
Huddle h1 & clock & minute 16 & pulled in by the clock\\
\bottomrule\end{tabular}
\subsection*{Targets}\begin{itemize}\item I can name the two objects in any force interaction.
\item I can sort forces as contact or field forces.
\item I can draw a free-body diagram and find the net force.\end{itemize}
\subsection*{Predictions}\begin{itemize}
\item A book rests on a table. How many forces act on the book? \A{Two: gravity down and the table's push up.} Gravity pulls down and the table pushes up. They are equal, so the net force is zero.
\item A crate slides to the right at constant speed. A person pushes it. How does the push compare with friction? \A{The push equals friction.} Constant speed means zero net force, so the push and friction are equal.
\item The same crate now speeds up to the right. Which way does the net force point? \A{To the right.} The net force points the way the crate speeds up: the push beats friction.
\end{itemize}
\section*{Huddle script (answers in green)}
\subsection*{Definitions}\begin{itemize}
\item \textbf{Free-body diagram:} A \A{picture} of one object with an \A{arrow} for each force on it.
\item \textbf{Normal force:} The push of a \A{surface} on an object, \A{perpendicular} to the surface.
\item \textbf{Net force:} The \A{sum} of all the forces on an object, with \A{direction}.
\end{itemize}
\subsection*{Worked example}
A crate has W = 40 N down, N = 40 N up, a push of 20 N right, and friction of 8 N left. Find the net force.\\
\begin{itemize}
\item Up and down: 40 N and 40 N \A{cancel}.
\item Right and left: 20 -- 8 = \A{12} N.
\item Net force: 12 N to the \A{right}.
\item The crate \A{speeds up} to the right.
\end{itemize}
\subsection*{Label the crate}
\begin{enumerate}
\item The \A{applied force} pushes right.
\item \A{Friction} opposes the slide.
\item \A{Gravity} pulls down.
\item The \A{normal force} pushes up.
\end{enumerate}
\section*{Every level, with answers}
\subsection*{Name the objects in an interaction}
\noindent\textbf{Level 1.} A person pushes a wall. Which object exerts the push on the wall?\\ \A{(A) The person} \quad (B) The wall \quad (C) The floor \quad (D) The air\\ \A{Answer: A} \quad The person's hands touch the wall and push it.\\ \textit{Misconception targeted: Students let an object exert a force on itself, or name "motion" as the source.}\par\medskip
\noindent\textbf{Level 2.} A book rests on a table. Which two objects interact through the normal force?\\ \A{(A) The book and the table} \quad (B) The book and the floor \quad (C) The book and the air \quad (D) The table and the floor\\ \A{Answer: A} \quad The normal force is a push between two surfaces that touch: the book and the table.\\ \textit{Misconception targeted: Students name one object, or a pair that does not touch.}\par\medskip
\noindent\textbf{Level 3.} Which force would still act on a satellite-free rock in empty space near the Earth?\\ \A{(A) Gravity} \quad (B) Normal force \quad (C) Friction \quad (D) Tension\\ \A{Answer: A} \quad Gravity is a field force: the Earth pulls without touching. Normal force, tension and friction need contact.\\ \textit{Misconception targeted: Students call gravity a contact force because it "touches the ground".}\par\medskip
\noindent\textbf{Level 4.} A banner hangs from a rope fixed to the ceiling. Which list names every object that exerts a force ON the banner?\\ \A{(A) The rope and the Earth} \quad (B) The rope, the Earth and the ceiling \quad (C) The rope only \quad (D) The ceiling and the Earth\\ \A{Answer: A} \quad The rope touches the banner and pulls up. The Earth pulls down at a distance. The ceiling touches the rope, not the banner.\\ \textit{Misconception targeted: Students include every object in the scene, not just those acting on this object.}\par\medskip
\subsection*{Sort and name the forces}
\noindent\textbf{Level 1.} A rough floor opposes a sliding box. This force is called...\\ \A{(A) Friction} \quad (B) Normal force \quad (C) Gravity \quad (D) Tension\\ \A{Answer: A} \quad A force that opposes sliding between two surfaces is friction.\\ \textit{Misconception targeted: Students mix up normal force and gravity because both are vertical.}\par\medskip
\noindent\textbf{Level 2.} A book rests on a flat table. Which way does the normal force on the book point?\\ \A{(A) Up} \quad (B) Down \quad (C) Left \quad (D) Right\\ \A{Answer: A} \quad The normal force is perpendicular to the surface and pushes away from it: up.\\ \textit{Misconception targeted: Students think friction points the way the object moves, or that gravity stops at the top of a throw.}\par\medskip
\noindent\textbf{Level 3.} Is applied force a contact force or a field force?\\ \A{(A) Contact force} \quad (B) Field force \quad (C) Neither\\ \A{Answer: A} \quad Applied force needs the two objects to touch: a contact force.\\ \textit{Misconception targeted: Students call gravity a contact force, or call normal force a field force.}\par\medskip
\noindent\textbf{Level 4.} A cart is pushed right across a rough floor. Which list names every force on it?\\ \A{(A) W, N, F, f} \quad (B) W, N, F \quad (C) W, N, F, f, force of motion \quad (D) F, f\\ \A{Answer: A} \quad Gravity (W) and the floor's normal force (N) act on it, plus the pull or push and friction from the rough floor. Nothing else touches it.\\ \textit{Misconception targeted: Students add a "force of motion", or forget the normal force and friction.}\par\medskip
\subsection*{Pick the free-body diagram}
\noindent\textbf{Level 1.} A phone lies on a shelf. Which diagram shows the forces on the phone?\\ \A{(A) W down, N up; N equals W} \quad (B) W down \quad (C) W down, N up, F right; N equals W \quad (D) W down, N up; N shorter than W\\ \A{Answer: A} \quad Gravity pulls down and the surface pulls or pushes up with an equal arrow, so the net force is zero.\\ \textit{Misconception targeted: Students draw only gravity on a still object, or add a force for "staying there".}\par\medskip
\noindent\textbf{Level 2.} A crate slides left at constant speed. A person pushes it left. Which diagram is right?\\ \A{(A) W down, N up, F left, f right; N equals W; F equals f} \quad (B) W down, N up, F left; N equals W \quad (C) W down, N up, F left, f right, M left; N equals W; F equals f \quad (D) W down, N up, F left, f right; N equals W; F longer than f\\ \A{Answer: A} \quad Constant speed means net force zero: push equals friction, and N equals W.\\ \textit{Misconception targeted: Students believe motion needs a bigger forward force than friction, or draw a "force of motion".}\par\medskip
\noindent\textbf{Level 3.} A trunk speeds up to the right. A person pushes it right. Which diagram is right?\\ \A{(A) W down, N up, F right, f left; N equals W; F longer than f} \quad (B) W down, N up, F right, f left; N equals W; F equals f \quad (C) W down, N up, F right, f left, M right; N equals W; F longer than f \quad (D) W down, N up, F right; N equals W\\ \A{Answer: A} \quad Speeding up means the net force points right: the push arrow is longer than friction. N still equals W.\\ \textit{Misconception targeted: Students draw push and friction equal for any motion.}\par\medskip
\noindent\textbf{Level 4.} A person pushes a crate right with 35 N. Friction is 10 N and the crate weighs 40 N. It speeds up. Which diagram has arrows drawn to scale?\\ \A{(A) W down (40 N), N up (40 N), F right (35 N), f left (10 N); N equals W; F longer than f} \quad (B) W down (40 N), N up (30 N), F right (35 N), f left (10 N); N shorter than W; F longer than f \quad (C) W down (40 N), N up (40 N), F right (10 N), f left (35 N); N equals W; F shorter than f \quad (D) W down (40 N), N up (40 N), F right (35 N); N equals W\\ \A{Answer: A} \quad N equals W (40 N). The right arrow (35 N) is longer than the left (10 N), so the net force is 25 N to the right.\\ \textit{Misconception targeted: Students draw lengths that do not match the numbers, or let N differ from W.}\par\medskip
\subsection*{Which way is the net force?}
\noindent\textbf{Level 1.} A box has 6 N pushing right and 6 N pushing left. Which way is the net force?\\ \A{(A) It is zero} \quad (B) Right \quad (C) Left\\ \A{Answer: A} \quad Equal arrows in opposite directions cancel: the net force is zero.\\ \textit{Misconception targeted: Students think equal forces still "push" the box one way.}\par\medskip
\noindent\textbf{Level 2.} A crate has 9 N pulling left and 7 N of friction pulling right. Which way is the net force?\\ \A{(A) Left} \quad (B) Right \quad (C) It is zero\\ \A{Answer: A} \quad The bigger force is 9 N, left. The net force is 2 N left.\\ \textit{Misconception targeted: Students take the net force along the motion, or call friction the "smaller force" and ignore it.}\par\medskip
\noindent\textbf{Level 3.} A crate weighs 60 N, the floor pushes up 60 N, a person pushes right 8 N, and friction is 4 N left. Which way is the net force?\\ \A{(A) Right} \quad (B) It is zero \quad (C) Left \quad (D) Up\\ \A{Answer: A} \quad Up and down cancel (60 N each). Right 8 N and left 4 N: net 4 N right.\\ \textit{Misconception targeted: Students try to add vertical and horizontal forces into one number.}\par\medskip
\noindent\textbf{Level 4.} A 3 kg lamp (weight 29.4 N) hangs from a chain pulling up 18 N. Which way is the net force?\\ \A{(A) Down} \quad (B) Up \quad (C) It is zero\\ \A{Answer: A} \quad W = 3 $\times$ 9.8 = 29.4 N. The chain pulls 18 N. The weight is bigger, so the net force points down.\\ \textit{Misconception targeted: Students assume a hanging object is always balanced.}\par\medskip
\subsection*{Find the net force}
\noindent\textbf{Level 1.} A cart is pushed right with 9 N and friction pulls left with 4 N. How big is the net force?\\ \A{Answer: 5 N} \quad 9 -- 4 = 5 N, to the right.\\ \textit{Misconception targeted: Students add forces that point opposite ways.}\par\medskip
\noindent\textbf{Level 2.} A crate is pulled right with 12.5 N. Friction is 8.5 N left. How big is the net force?\\ \A{Answer: 4 N} \quad 12.5 -- 8.5 = 4 N.\\ \textit{Misconception targeted: Students drop the decimal or add the forces.}\par\medskip
\noindent\textbf{Level 3.} A 5 kg box (weight 49 N, floor pushes up 49 N) is pushed right with 15 N. Friction is 10 N. What is the net force?\\ \A{Answer: 5 N} \quad Up and down cancel. Net = 15 -- 10 = 5 N, to the right.\\ \textit{Misconception targeted: Students include weight in a horizontal net force.}\par\medskip
\noindent\textbf{Level 4.} Two people push a crate right with 23 N and 14 N while friction pulls left with 15 N. What is the net force?\\ \A{Answer: 22 N} \quad Right: 23 + 14 = 37 N. Left: 15 N. Net = 37 -- 15 = 22 N right.\\ \textit{Misconception targeted: Students subtract in the order they read the numbers.}\par\medskip
\subsection*{Two objects, one diagram}
\noindent\textbf{Level 1.} A cup sits on a desk. Which force does NOT act ON the cup?\\ \A{(A) The cup's push on the desk} \quad (B) Gravity on the cup \quad (C) The desk's push up on the cup\\ \A{Answer: A} \quad The cup pushes on the desk, so that force acts on the desk. A diagram of the cup shows only forces on the cup.\\ \textit{Misconception targeted: Students draw the object's force on the table as a force on the object.}\par\medskip
\noindent\textbf{Level 2.} A book A sits on a book B, and B sits on a table. Which list names the forces on book A?\\ \A{(A) Gravity on A, and B's push up} \quad (B) Gravity on A, B's push up, and the table's push up \quad (C) Gravity on A and on B \quad (D) Gravity on A only\\ \A{Answer: A} \quad A is touched only by B, so B pushes up on it. The table touches B, not A.\\ \textit{Misconception targeted: Students list every force in the picture on every object.}\par\medskip
\noindent\textbf{Level 3.} Box A (4 kg) sits on box B (1 kg). The table pushes up on B. How hard does the table push, in newtons? (g = 9.8 N/kg)\\ \A{Answer: 49 N} \quad The table holds up both: N = (4 + 1) $\times$ 9.8 = 49 N.\\ \textit{Misconception targeted: Students count only the lower box's weight.}\par\medskip
\noindent\textbf{Level 4.} Boxes of 2 kg and 3 kg are tied together and pulled right with 34 N. Friction on the pair totals 12 N. What is the net force on the pair?\\ \A{Answer: 22 N} \quad 34 -- 12 = 22 N to the right. The tension between the boxes is inside the system and does not count.\\ \textit{Misconception targeted: Students add the tension between the boxes into the net force.}\par\medskip
\section*{Common misconceptions}\begin{tabularx}{\textwidth}{@{}X X@{}}\toprule Students say & Say instead\\\midrule
The crate needs a force of motion pushing it forward. & Motion is not a force. Name the object that exerts every arrow. At constant speed the net force is zero.\\
A book at rest has no forces on it, or only gravity. & The table pushes up. At rest the forces balance: N equals W.\\
The ceiling pulls on the hanging lamp. & Only objects touching the lamp (the chain) and the Earth act on it.\\
Gravity is a contact force. & Gravity is a field force. It acts without touching.\\
Friction points the way the object moves. & Friction opposes the sliding between surfaces.\\
Add all the numbers to get the net force. & Add forces that point the same way. Subtract opposite ones. Vertical and horizontal do not mix.\\
\bottomrule\end{tabularx}
\end{document}